%% Dans ce fichier, les �ventuels accents sont cod�s
%% en ISO-8859-15 (aka latin9)
%% La ligne suivante donne trois exemples 
%% e accent aigu : �, e accent grave : �, e-dans-l'o : �
%%%%%ENTREES%%%%
\begin{tabarg}
\multicolumn{3}{l}{\titretab {Entr\'ees}}\\
\hline
$X$ & Vecteur, \texttt{float} & Param�tre d'entr�e (unit� quelconque)\\
$Seuil$ & Coeff., \texttt{float} & Seuil pour le choix du polyn\^ome � appliquer \\
$a_0$ & Coeff., \texttt{float} & Coefficient du terme de degr� 0 du 1er polyn\^ome\\
$a_1$ & Coeff., \texttt{float} & Coefficient du terme de degr� 1 du 1er polyn\^ome\\
$b_0$ & Coeff., \texttt{float} & Coefficient du terme de degr� 0 du 2�me polyn\^ome\\
$b_1$ & Coeff., \texttt{float} & Coefficient du terme de degr� 1 du 2�me polyn\^ome\\

$Unit$ & Coeff., \texttt{String} & Unit� du param�tre r�sultat\\
\end{tabarg}

\medskip
%%%%%SORTIES%%%%
\begin{tabarg}
\multicolumn{3}{l}{\titretab {Sorties}}\\
\hline
$Y$ & Vecteur, \texttt{float} & Param�tre r�sultat ($Unit$)\\
\end{tabarg}
